% Gaussian elimination / Cholesky fill-in when eliminating an interior vertex i.
\begin{tikzpicture}[x=1mm,y=1mm,font=\figurefont\fontsize{8}{10}\selectfont,>=Stealth,
  vertex/.style={circle,draw=FigureBlue,fill=FigureBlueFill,inner sep=0pt,
    minimum size=5mm,line width=.55pt},
  elim/.style={circle,draw=FigureYellow,fill=FigureYellowFill,inner sep=0pt,
    minimum size=5.2mm,line width=.65pt},
  edge/.style={draw=FigureStroke,line width=.65pt},
  fillin/.style={draw=FigureRed,dashed,line width=.8pt},
  note/.style={font=\FigureNoteFont\fontsize{7}{8.5}\selectfont,text=FigureInk,inner sep=1pt},
  panel/.style={font=\figurefont\fontsize{7.5}{9}\selectfont,text=FigureInk}]
  \definecolor{FigureBlue}{HTML}{7188A3}
  \definecolor{FigureBlueFill}{HTML}{E5EAF0}
  \definecolor{FigureGreen}{HTML}{829C8B}
  \definecolor{FigureGreenFill}{HTML}{E8EEE9}
  \definecolor{FigureYellow}{HTML}{B4A16F}
  \definecolor{FigureYellowFill}{HTML}{F1ECDD}
  \definecolor{FigureRed}{HTML}{A7787D}
  \definecolor{FigureInk}{HTML}{3B4252}
  \definecolor{FigureMuted}{HTML}{6B7280}
  \definecolor{FigureStroke}{HTML}{4C4D4F}

  % Panel (a): Before eliminating vertex i
  \begin{scope}[xshift=0mm]
    \node[elim] (i) at (14,9) {$i$};
    \node[vertex] (j) at (14,21) {$j$};
    \node[vertex] (k) at (2,1) {$k$};
    \node[vertex] (l) at (26,1) {$\ell$};
    \draw[edge] (i)--(j) node[midway,right,note] {$c_{i,j}$};
    \draw[edge] (i)--(k) node[midway,above left,note] {$c_{i,k}$};
    \draw[edge] (i)--(l) node[midway,above right,note] {$c_{i,\ell}$};
    \node[panel] at (14,-7) {(a) 消去 $v_i$ 前的局部星形连接};
  \end{scope}

  \draw[->,FigureStroke,line width=.75pt] (33,9) -- node[above,note] {消去 $v_i$} (47,9);

  % Panel (b): After eliminating vertex i (clique fill-in on neighbors j, k, l)
  \begin{scope}[xshift=54mm]
    \node[vertex] (bj) at (14,21) {$j$};
    \node[vertex] (bk) at (2,1) {$k$};
    \node[vertex] (bl) at (26,1) {$\ell$};
    \draw[fillin] (bj)--(bk);
    \draw[fillin] (bj)--(bl);
    \draw[fillin] (bk)--(bl);
    \node[note,text=FigureRed] at (14,8) {新增填充边};
    \node[panel] at (14,-7) {(b) 消元后邻居间产生填充耦合};
  \end{scope}
\end{tikzpicture}
